\subsection{The homogeneous positive-rational index}
\label{sec:qca_gnvw_index}

Let $d>0$ and let $\omega$ be a nearest-neighbour star-automorphism of the
homogeneous chain. The following uses the fixed nearest-neighbour cell
presentation. The site-dependent source setting is recorded separately in
\cite{gap:gnvw12_site_dependent_adjacent_generation_scope}.

\begin{definition}[The fixed-cut GNVW index]
    \label{def:qca_gnvw_index}
    \lean{SpinChain.PropagatesWithin.supportAlgebraMatrixSize,
        SpinChain.PropagatesWithin.supportAlgebraMatrixSize_spec,
        SpinChain.PropagatesWithin.supportAlgebraMatrixSize_neZero,
        SpinChain.PropagatesWithin.gnvwIndex}
    \leanok
    \uses{lem:qca_physical_pair_adjacent_generation}
    Choose the coherent positive matrix-size family $r(y)$ of the actual
    support factors. Define
    \begin{align}
        \operatorname{ind}(\omega) & =\frac{r(0)}{d}\in\mathbb Q_{>0}.
        \label{eq:qca_gnvw_index_fixed_cut}
    \end{align}
    Its codomain is the multiplicative group of positive rationals.
    This is the homogeneous specialization of equation~dimtransfer in
    \cite[lines~1317--1329]{Gross2012IndexTheory}.
\end{definition}

\begin{lemma}[All cut presentations of the same index]
    \label{thm:qca_gnvw_index_dimtransfer}
    \lean{SpinChain.PropagatesWithin.supportAlgebraMatrixSize_even_add_one,
        SpinChain.PropagatesWithin.supportAlgebraMatrixSize_even_eq_zero,
        SpinChain.PropagatesWithin.coe_gnvwIndex_eq_even,
        SpinChain.PropagatesWithin.coe_gnvwIndex_eq_odd,
        SpinChain.PropagatesWithin.coe_gnvwIndex_eq_next_even,
        SpinChain.PropagatesWithin.gnvwIndex_ratio_eq,
        SpinChain.PropagatesWithin.coe_gnvwIndex_eq_of_evenSupportAlgebra_equiv,
        SpinChain.PropagatesWithin.supportAlgebraMatrixSize_even_eq_of_equiv,
        SpinChain.PropagatesWithin.supportAlgebraMatrixSize_odd_eq_of_equiv}
    \leanok
    \uses{def:qca_gnvw_index}
    For every $x\in\mathbb Z$,
    \begin{align}
        \operatorname{ind}(\omega)
         & =\frac{r(2x)}{d}
        =\frac{d}{r(2x+1)}
        =\frac{r(2x+2)}{d}.
        \label{eq:qca_gnvw_index_dimtransfer}
    \end{align}
    These are presentations using the same ordered support family. No
    translation covariance is assumed. The value is independent of the
    integer cut and of the chosen full-matrix presentation at that cut.
    This is homogeneous equation~dimtransfer of
    \cite[lines~1317--1324]{Gross2012IndexTheory}.
\end{lemma}
\begin{proof}\leanok
    The two size products equal $d^2$. Cancelling the positive common odd
    size gives $r(2x)=r(2x+2)$. Integer induction in both directions makes
    the even size constant. Dividing the first product identity gives
    the reciprocal odd ratio. Matrix-algebra isomorphisms identify the
    sizes of any two presentations of the same actual factor.
\end{proof}

\begin{lemma}[Support factors of a full tensor factor]
    \label{lem:qca_kronecker_embedding_supports}
    \lean{Matrix.leftSupportAlgebra_range_leftKroneckerEmbed,
        Matrix.rightSupportAlgebra_range_leftKroneckerEmbed,
        Matrix.leftSupportAlgebra_range_rightKroneckerEmbed,
        Matrix.rightSupportAlgebra_range_rightKroneckerEmbed}
    \leanok
    \uses{def:qca_left_support_algebra,def:qca_right_support_algebra}
    Let $I$ and $J$ be finite index sets. Then
    \begin{align}
        \operatorname{Spp}_{\mathrm L}
        \bigl(M_I(\mathbb C)\otimes I_J\bigr)
                                       & =M_I(\mathbb C)
        \quad\text{if }J\neq\emptyset, &
        \operatorname{Spp}_{\mathrm R}
        \bigl(M_I(\mathbb C)\otimes I_J\bigr)
                                       & =\mathbb C I_J,\notag   \\
        \operatorname{Spp}_{\mathrm L}
        \bigl(I_I\otimes M_J(\mathbb C)\bigr)
                                       & =\mathbb C I_I,       &
        \operatorname{Spp}_{\mathrm R}
        \bigl(I_I\otimes M_J(\mathbb C)\bigr)
                                       & =M_J(\mathbb C)
        \quad\text{if }I\neq\emptyset.
        \label{eq:qca_kronecker_embedding_supports}
    \end{align}
    Only the full support of each factor needs the spectator index set to
    be nonempty. These are the support calculations in the shift example
    following equation~RR2x of
    \cite[lines~1251--1268]{Gross2012IndexTheory}.
\end{lemma}
\begin{proof}\leanok
    For $A\in M_I(\mathbb C)$, $B\in M_J(\mathbb C)$, $i,j\in I$ and
    $r,s\in J$, the coefficients are
    \begin{align}
        (A\otimes I_J)^{\mathrm L}_{rs}
         & =\delta_{rs}A,    &
        (A\otimes I_J)^{\mathrm R}_{ij}
         & =A_{ij}I_J,\notag   \\
        (I_I\otimes B)^{\mathrm L}_{rs}
         & =B_{rs}I_I,       &
        (I_I\otimes B)^{\mathrm R}_{ij}
         & =\delta_{ij}B.
        \label{eq:qca_kronecker_embedding_coefficients}
    \end{align}
    Taking $r=s$, which needs $J\neq\emptyset$, recovers every $A$ as a
    left coefficient; taking $i=j$, which needs $I\neq\emptyset$, recovers
    every $B$ as a right coefficient. The remaining coefficients are scalar
    multiples of the identity, so they generate only the scalars.
\end{proof}

\begin{lemma}[Finite-region configuration count]
    \label{lem:qca_finite_region_configuration_count}
    \lean{SpinChain.card_config}
    \leanok
    A region $\Lambda$ has $d^{|\Lambda|}$ configurations. In particular,
    a physical pair has Hilbert-space dimension $d^2$ and observable-algebra
    dimension $d^4$.
\end{lemma}
\begin{proof}\leanok
    There is one independent choice among $d$ values at each site.
\end{proof}

\begin{lemma}[Faithful bipartite observable coordinates]
    \label{lem:qca_bipartite_observable_embedding}
    \lean{SpinChain.bipartiteMatrixToQuasiLocalObservable,
        SpinChain.bipartiteMatrixToQuasiLocalObservable_injective,
        SpinChain.bipartiteMatrixToQuasiLocalObservable_leftKroneckerEmbed,
        SpinChain.bipartiteMatrixToQuasiLocalObservable_rightKroneckerEmbed,
        SpinChain.map_range_leftKroneckerEmbed_bipartiteMatrixToQuasiLocalObservable,
        SpinChain.map_range_rightKroneckerEmbed_bipartiteMatrixToQuasiLocalObservable}
    \leanok
    For disjoint finite regions, the canonical bipartite matrix coordinates
    embed faithfully into the quasi-local algebra. The full left and right
    factor ranges become exactly the observable algebras of their respective
    finite regions. These are the finite coordinates of
    \cite[lines~1251--1278]{Gross2012IndexTheory}.
\end{lemma}
\begin{proof}\leanok
    Write $\varphi$ for the bipartite embedding and
    $\iota_\Lambda,\iota_\Gamma$ for the canonical finite-region
    observable inclusions. The ordered bipartite coordinate equivalence
    gives, for $A\in M_{\Lambda}(\mathbb C)$ and
    $B\in M_{\Gamma}(\mathbb C)$,
    \begin{align}
        \varphi(A\otimes I_\Gamma)
         & =\iota_\Lambda(A), &
        \varphi(I_\Lambda\otimes B)
         & =\iota_\Gamma(B).
        \label{eq:qca_bipartite_observable_embedding}
    \end{align}
    Injectivity of $\varphi$ follows from faithfulness of the canonical
    inclusion of $\Lambda\cup\Gamma$. The two identities map the full left
    and right factor ranges onto the ranges of $\iota_\Lambda$ and
    $\iota_\Gamma$.
\end{proof}

\begin{lemma}[Direct shift normalization]
    \label{thm:qca_gnvw_index_shift}
    \lean{SpinChain.quasiLocalTranslation_propagatesWithin_singleton,
        SpinChain.quasiLocalTranslation_propagatesWithin_neg_one,
        SpinChain.quasiLocalTranslation_propagatesWithin_one,
        SpinChain.range_quasiLocalTranslation_comp_quasiLocalObservable,
        SpinChain.translateRegion_neg_one_evenPair,
        SpinChain.translateRegion_one_evenPair,
        SpinChain.PropagatesWithin.evenPairLocalImage_quasiLocalTranslation_neg_one,
        SpinChain.PropagatesWithin.evenPairLocalImage_quasiLocalTranslation_one,
        SpinChain.PropagatesWithin.supportAlgebras_quasiLocalTranslation_neg_one,
        SpinChain.PropagatesWithin.supportAlgebras_quasiLocalTranslation_one,
        SpinChain.PropagatesWithin.supportAlgebraMatrixSize_quasiLocalTranslation_neg_one_even,
        SpinChain.PropagatesWithin.supportAlgebraMatrixSize_quasiLocalTranslation_neg_one_odd,
        SpinChain.PropagatesWithin.supportAlgebraMatrixSize_quasiLocalTranslation_one_even,
        SpinChain.PropagatesWithin.supportAlgebraMatrixSize_quasiLocalTranslation_one_odd,
        SpinChain.PropagatesWithin.gnvwIndex_quasiLocalTranslation_neg_one,
        SpinChain.PropagatesWithin.gnvwIndex_quasiLocalTranslation_one}
    \leanok
    \uses{def:qca_gnvw_index,thm:qca_gnvw_index_dimtransfer,
        lem:qca_kronecker_embedding_supports,lem:qca_bipartite_observable_embedding,
        lem:qca_finite_region_configuration_count}
    Write $T_a$ for observable translation by $a$. The source Heisenberg
    shift is $T_{-1}$, which moves observables left and states right.
    Directly in the ordered finite coordinates,
    \begin{align}
        \mathcal R_{2x}(T_{-1})   & =\mathcal A_{L_x},  &
        \mathcal R_{2x+1}(T_{-1}) & =\mathbb C I,\notag   \\
        \mathcal R_{2x}(T_1)      & =\mathbb C I,       &
        \mathcal R_{2x+1}(T_1)    & =\mathcal A_{P_x}.
        \label{eq:qca_translation_support_factors}
    \end{align}
    Hence
    \begin{align}
        \operatorname{ind}(T_{-1}) & =d,      &
        \operatorname{ind}(T_1)    & =d^{-1}.
        \label{eq:qca_gnvw_index_shift}
    \end{align}
    The first is the homogeneous shift normalization of
    \cite[lines~1325--1329]{Gross2012IndexTheory}; the second is a direct
    orientation check. The repeated even subscript in the source scalar
    support example at line~1268 is corrected as documented in
    \cite{gap:gnvw12_shift_example_subscript}.
\end{lemma}
\begin{proof}\leanok
    Translation has its singleton propagation neighbourhood, which lies
    inside the nearest-neighbour interval. It maps the full source-pair
    algebra onto the left pair for $T_{-1}$, and onto the right pair for
    $T_1$. Faithful bipartite coordinates identify the finite images with
    the full left and right tensor factors respectively. Their supports
    therefore have the displayed full/scalar values. The full factor has
    matrix size $d^2$; the companion product identity makes the scalar
    factor's matrix size one. Substitution into the already derived index
    ratios proves both values without any composition or inversion law.
\end{proof}
